\Needspace{5\baselineskip}\textbf{Problems}

\begin{enumerate}[label=\thesection.\arabic*.,ref=\thesection.\arabic*,start=24,series=book-problems]

\item\label{ex:complex-matrix-banach-hilbert-adjoints} 设 $A$ 为 $m\times n$ 复矩阵, 并通过矩阵乘法定义从 $\mathbb{C}^n\to\mathbb{C}^m$ 的算子 (也记为 $A$). 算子 $A$ 的伴随, Hilbert 伴随以及矩阵 $A$ 之间有何关系?

\item\label{ex:hilbert-adjoint-product-reversal} 若 $A$ 和 $B$ 属于 $\mathcal{L}(H)$, 证明\[(AB)^*=B^*A^*.\]

\item\label{ex:operator-extension-adjoint-reverse-extension} 证明, 若 $(\mathcal{D}(B),B)$ 是 $(\mathcal{D}(A),A)$ 的扩张, 则 $(\mathcal{D}(A^\times),A^\times)$ 是 $(\mathcal{D}(B^\times),B^\times)$ 的扩张.

\item\label{ex:adjoint-closure-theorem-proof-completion} 完成 Theorem~\ref{thm:adjoint-closedness-and-double-adjoint-closure} 的证明.

\item\label{ex:right-shift-hilbert-adjoint} 计算 Example~\ref{ex:right-left-shifts-domain-dependent-inverses} 中定义的右移算子 $S_r$ 的 Hilbert 伴随.

\item\label{ex:second-derivative-symmetric-not-self-adjoint} 设 $H=L^2(0,1)$, 并令\[\mathcal{D}(A)=\{u\in H^2(0,1)\mid u(0)=u'(0)=u(1)=0\}.\]这里的边界条件在迹的意义下理解. 用下式定义 $A:\mathcal{D}(A)\to H$:\[A(u):=-\dfrac{\mathrm{d}^2u}{\mathrm{d}x^2}.\]求 $(\mathcal{D}(A),A)$ 的 Hilbert 伴随. 该算子是否对称, 是否自伴随?

\item\label{ex:classical-laplacian-boundary-domains-symmetric} 证明 Example~\ref{ex:laplacian-dirichlet-neumann-domains} 中定义的 $(\widetilde{\mathcal{D}}_D(\Delta),\Delta)$ 和 $(\widetilde{\mathcal{D}}_N(\Delta),\Delta)$ 是对称的.

\item\label{ex:adjoint-deficiency-eigenvalue-proof} 证明 Theorem~\ref{thm:adjoint-deficiency-conjugate-eigenvalue}.

\item\label{ex:hilbert-adjoint-product-norm-square} 设 $A\in\mathcal{L}(H)$. 证明 $\|A^*A\|=\|A\|^2$.

\item\label{ex:self-adjoint-spectral-radius-equals-norm} 可以证明, 对于算子 $A\in\mathcal{L}(X)$, 有\[r_\sigma(A)=\lim_{n\to\infty}\|A^n\|^{1/n}.\]利用这一事实和 Problem~\ref{ex:hilbert-adjoint-product-norm-square}, 证明若 $A\in\mathcal{L}(H)$ 是自伴随的, 则 $r_\sigma(A)=\|A\|$.

\item\label{ex:commuting-operator-product-spectral-radius} 假设 $A,B\in\mathcal{L}(X)$ 且 $AB=BA$. 证明\[r_\sigma(AB)\leq r_\sigma(A)r_\sigma(B).\]证明这一结果中的交换性假设不可缺少.

\item\label{ex:symmetric-operator-closable} 证明每个对称算子都是可闭的.

\item\label{ex:hilbert-adjoint-bounded-inverse-equivalence} 设 $X$ 和 $Y$ 为 Hilbert 空间, 并假设 $A\in\mathcal{L}(X,Y)$ 是双射. 证明, $A$ 有有界逆, 当且仅当 $A^*$ 有有界逆.

\item\label{ex:essential-self-adjointness-deficiency-proof} 证明 Theorem~\ref{thm:essential-self-adjointness-zero-deficiency}.

\item\label{ex:right-left-shift-spectra} 描述 Example~\ref{ex:right-left-shifts-domain-dependent-inverses} 中所述右移算子和左移算子的谱.

\end{enumerate}
